% PDF page 26
\source{gilbarg-trudinger:page-0026}



\begin{theorem}\label{gt.thm2.1}\dcref{Theorem 2.1}\group{chapter-2}\source{gilbarg-trudinger:page-0026}
\emph{Let $u \in C^2(\Omega)$ satisfy $\Delta u = 0$ $(\geq 0,\ \leq 0)$ in $\Omega$. Then for any ball
$B = B_R(y) \subset\subset \Omega$, we have}

\begin{equation}
u(y) = (\leq,\ \geq)\ \frac{1}{n\omega_n R^{n-1}} \int_{\partial B} u\,ds,
\tag{2.5}
\end{equation}

\begin{equation}
u(y) = (\leq,\ \geq)\ \frac{1}{\omega_n R^n} \int_{B} u\,dx.
\tag{2.6}
\end{equation}

\noindent For harmonic functions, Theorem 2.1 thus asserts that the function value at the center of the ball $B$ is equal to the integral mean values over both the surface $\partial B$ and $B$ itself. These results, known as the \emph{mean value theorems}, in fact also characterize harmonic functions; (see Theorem 2.7).

\medskip

\begin{proof}
Let $\rho \in (0, R)$ and apply the identity (2.4) to the ball $B_\rho = B_\rho(y)$. We obtain

\begin{equation*}
\int_{\partial B_\rho} \frac{\partial u}{\partial \nu}\,ds = \int_{B_\rho} \Delta u\,dx = (\geq,\ \leq) 0.
\end{equation*}

\noindent Introducing radial and angular coordinates $r = |x-y|$, $\omega = \dfrac{x-y}{r}$, and writing
$u(x) = u(y + r\omega)$, we have

\begin{align*}
\int_{\partial B_\rho} \frac{\partial u}{\partial \nu}\,ds
&= \int_{\partial B_\rho} \frac{\partial u}{\partial r}(y+\rho\omega)\,ds
 = \rho^{n-1}\int_{|\omega|=1}\frac{\partial u}{\partial r}(y+\rho\omega)\,d\omega \\
&= \rho^{n-1}\frac{\partial}{\partial \rho}\int_{|\omega|=1} u(y+\rho\omega)\,d\omega
 = \rho^{n-1}\frac{\partial}{\partial \rho}\left[\rho^{1-n}\int_{\partial B_\rho} u\,ds\right] \\
&= (\geq,\ \leq) 0.
\end{align*}

\noindent Consequently for any $\rho \in (0, R)$,

\begin{equation*}
\rho^{1-n}\int_{\partial B_\rho} u\,ds = (\leq,\ \geq) R^{1-n}\int_{\partial B_R} u\,ds
\end{equation*}

\noindent and since

\begin{equation*}
\lim_{\rho \to 0} \rho^{1-n}\int_{\partial B_\rho} u\,ds = n\omega_n u(y)
\end{equation*}

\noindent relations (2.5) follow. To get the solid mean value inequalities, that is, relations (2.6), we write (2.5) in the form

\begin{equation*}
n\omega_n \rho^{n-1}u(y) = (\leq,\ \geq)\int_{\partial B_\rho} u\,ds,\qquad \rho \leq R
\end{equation*}

\noindent and integrate with respect to $\rho$ from $0$ to $R$. The relations (2.6) follow immediately.\hfill$\square$
\end{proof}
\end{theorem}


% PDF page 27
\source{gilbarg-trudinger:page-0027}

\noindent\textbf{2.2. Maximum and Minimum Principle}

\medskip

\noindent\textbf{2.2. Maximum and Minimum Principle}

\medskip

With the aid of Theorem 2.1 the \emph{strong maximum principle} for subharmonic functions and the \emph{strong minimum principle} for superharmonic functions may be derived.

\medskip

\begin{theorem}\label{gt.thm2.2}\dcref{Theorem 2.2}\group{chapter-2}\source{gilbarg-trudinger:page-0027}
\emph{Let $\Delta u \geq 0$ $(\leq 0)$ in $\Omega$ and suppose there exists a point $y \in \Omega$ for which}
\[
u(y)=\sup_{\Omega} u \ (\inf_{\Omega} u).
\]
\emph{Then $u$ is constant. Consequently a harmonic function cannot assume an interior maximum or minimum value unless it is constant.}

\medskip

\begin{proof}\uses{gt.thm2.1}
Let $\Delta u \geq 0$ in $\Omega$, $M=\sup_{\Omega} u$ and define $\Omega_M=\{x\in \Omega \mid u(x)=M\}$. By assumption $\Omega_M$ is not empty. Furthermore since $u$ is continuous, $\Omega_M$ is closed relative to $\Omega$. Let $z$ be any point in $\Omega_M$ and apply the mean value inequality (2.6) to the subharmonic function $u-M$ in a ball $B=B_R(z)\subset\Omega$. We therefore obtain
\[
0=u(z)-M\leq \frac{1}{\omega_n R^n}\int_B (u-M)\, dx\leq 0,
\]
so that $u=M$ in $B_R(z)$. Consequently $\Omega_M$ is also open relative to $\Omega$. Hence $\Omega_M=\Omega$.
The result for superharmonic functions follows by replacement of $u$ by $-u$. \hfill $\square$
\end{proof}
\end{theorem}

\medskip

The strong maximum and minimum principles immediately imply global estimates, namely the following \emph{weak maximum and minimum principles}.

\medskip

\begin{theorem}\label{gt.thm2.3}\dcref{Theorem 2.3}\group{chapter-2}\source{gilbarg-trudinger:page-0027}
\emph{Let $u \in C^2(\Omega)\cap C^0(\overline{\Omega})$ with $\Delta u \geq 0$ $(\leq 0)$ in $\Omega$. Then, provided $\Omega$ is bounded,}
\begin{equation}
\tag{2.7}
\sup_{\Omega} u=\sup_{\partial\Omega} u \quad \left(\inf_{\Omega} u=\inf_{\partial\Omega} u\right).
\end{equation}

\medskip

\noindent\emph{Consequently, for harmonic $u$}
\[
\inf_{\partial\Omega} u\leq u(x)\leq \sup_{\partial\Omega} u, \qquad x\in\Omega.
\]
\end{theorem}

\medskip

A uniqueness theorem for the classical Dirichlet problem for Laplace's and Poisson's equation in bounded domains now follows from Theorem 2.3.

\medskip

\begin{theorem}\label{gt.thm2.4}\dcref{Theorem 2.4}\group{chapter-2}\source{gilbarg-trudinger:page-0027}
\emph{Let $u, v \in C^2(\Omega)\cap C^0(\overline{\Omega})$ satisfy $\Delta u=\Delta v$ in $\Omega$, $u=v$ on $\partial\Omega$. Then $u=v$ in $\Omega$.}

\medskip

\begin{proof}\uses{gt.thm2.3}
Let $w=u-v$. Then $\Delta w=0$ in $\Omega$ and $w=0$ on $\partial\Omega$. It follows from Theorem 2.3 that $w=0$ in $\Omega$. \hfill $\square$
\end{proof}
\end{theorem}

\medskip

\noindent Note that also by Theorem 2.3, we have that if $u$ and $v$ are harmonic and subharmonic functions respectively, agreeing on the boundary $\partial\Omega$, then $v\leq u$ in $\Omega$.


% PDF page 28
\source{gilbarg-trudinger:page-0028}


\noindent Hence the term subharmonic. A corresponding remark is true for superharmonic
functions. Later in this chapter, we employ this property of $C^2(\Omega)$ subharmonic
and superharmonic functions to expand their definition to larger classes of func-
tions. In the next chapter, an alternate method of proof for Theorems 2.2, 2.3 and
2.4 will be supplied when we treat maximum principles for general elliptic equa-
tions; (see also Problem 2.1).

\bigskip

\noindent\textbf{2.3.\ The Harnack Inequality}

\medskip

\noindent A further consequence of Theorem 2.1 is the following Harnack inequality for
harmonic functions.

\medskip

\begin{theorem}\label{gt.thm2.5}\dcref{Theorem 2.5}\group{chapter-2}\source{gilbarg-trudinger:page-0028}
\emph{Let $u$ be a non-negative harmonic function in $\Omega$. Then for any bounded
subdomain $\Omega' \subset\subset \Omega$, there exists a constant $C$ depending only on $n$, $\Omega'$ and $\Omega$ such
that}

\begin{equation}
\sup_{\Omega'} u \leq C \inf_{\Omega'} u.
\tag{2.8}
\end{equation}

\begin{proof}\uses{gt.thm2.1}
Let $y \in \Omega$, $B_{4R}(y) \subset \Omega$. Then for any two points $x_1, x_2 \in B_R(y)$, we have
by the inequalities (2.6)

\begin{align*}
u(x_1) &= \frac{1}{\omega_n R^n} \int_{B_R(x_1)} u\,dx \leq \frac{1}{\omega_n R^n} \int_{B_{2R}(y)} u\,dx, \\
u(x_2) &= \frac{1}{\omega_n (3R)^n} \int_{B_{3R}(x_2)} u\,dx \geq \frac{1}{\omega_n (3R)^n} \int_{B_{2R}(y)} u\,dx.
\end{align*}

\noindent Consequently we obtain

\begin{equation}
u(x_1) \leq 3^n u(x_2).
\tag{2.9}
\end{equation}

\noindent Now let $\Omega' \subset\subset \Omega$ and choose $x_1, x_2 \in \overline{\Omega}'$ so that $u(x_1)=\sup_{\Omega'} u$, $u(x_2)=\inf_{\Omega'} u$.
Let $\Gamma \subset \overline{\Omega}'$ be a closed arc joining $x_1$ and $x_2$, and choose $R$ so that $4R < \operatorname{dist}(\Gamma, \partial \Omega)$.
By virtue of the Heine-Borel theorem, $\Gamma$ can be covered by a finite number $N$
(depending only on $\Omega'$ and $\Omega$) of balls of radius $R$. Applying the estimate (2.9) in
each ball and combining the resulting inequalities, we obtain

\begin{equation*}
u(x_1) \leq 3^{nN} u(x_2).
\end{equation*}

\noindent Hence the estimate (2.8) holds with $C = 3^{nN}$. \hfill$\square$

\medskip

\noindent Note that the constant in (2.8) is invariant under similarity and orthogonal
transformations. A Harnack inequality for weak solutions of homogeneous elliptic
equations will be established in Chapter 8.
\end{proof}
\end{theorem}


% PDF page 29
\source{gilbarg-trudinger:page-0029}



\section*{2.4. Green's Representation}

As a prelude to existence considerations we derive now some further consequences
of the divergence theorem, namely, the Green identities. Let $\Omega$ be a domain for
which the divergence theorem holds and let $u$ and $v$ be $C^2(\overline{\Omega})$ functions. We select
$w = vDu$ in the identity (2.3) to obtain Green's first identity:

\begin{equation}
\int_{\Omega} v\,\Delta u\,dx + \int_{\Omega} Du \cdot Dv\,dx = \int_{\partial \Omega} v\,\frac{\partial u}{\partial \nu}\,ds.
\tag{2.10}
\end{equation}

Interchanging $u$ and $v$ in (2.10) and subtracting, we obtain Green's second identity:

\begin{equation}
\int_{\Omega} (v\,\Delta u-u\,\Delta v)\,dx = \int_{\partial \Omega} \left(v\,\frac{\partial u}{\partial \nu}-u\,\frac{\partial v}{\partial \nu}\right)\,ds.
\tag{2.11}
\end{equation}

Laplace's equation has the radially symmetric solution $r^{2-n}$ for $n>2$ and $\log r$ for
$n=2$, $r$ being radial distance from a fixed point. To proceed further from (2.11),
we fix a point $y$ in $\Omega$ and introduce the normalized fundamental solution of Laplace's
equation:

\begin{equation}
\Gamma(x-y)=\Gamma(|x-y|)=
\begin{cases}
\dfrac{1}{n(2-n)\omega_n}|x-y|^{2-n}, & n>2,\\[1.2ex]
\dfrac{1}{2\pi}\log |x-y|, & n=2.
\end{cases}
\tag{2.12}
\end{equation}

By simple computation we have

\[
D_i\Gamma(x-y)=\dfrac{1}{n\omega_n}(x_i-y_i)|x-y|^{-n};
\]

\begin{equation}
D_{ij}\Gamma(x-y)=\dfrac{1}{n\omega_n}\left\{|x-y|^2\delta_{ij}-n(x_i-y_i)(x_j-y_j)\right\}|x-y|^{-n-2}.
\tag{2.13}
\end{equation}

Clearly $\Gamma$ is harmonic for $x\neq y$. For later purposes we note the following derivative
estimates:

\begin{equation}
|D_i\Gamma(x-y)|\leq \dfrac{1}{n\omega_n}|x-y|^{1-n};
\end{equation}

\begin{equation}
|D_{ij}\Gamma(x-y)|\leq \dfrac{1}{\omega_n}|x-y|^{-n}.
\tag{2.14}
\end{equation}

\[
|D^\beta \Gamma(x-y)|\leq C|x-y|^{2-n-|\beta|},\qquad C=C(n,|\beta|).
\]

The singularity at $x=y$ prevents us from using $\Gamma$ in place of $v$ in Green's second
identity (2.11). One way of overcoming this difficulty is to replace $\Omega$ by $\Omega-\overline{B}_\rho$


% PDF page 30
\source{gilbarg-trudinger:page-0030}


where $B = B_\rho(y)$ for sufficiently small $\rho$. We can then conclude from (2.11) that

\begin{equation}
\tag{2.15}
\int_{\Omega-B_\rho} \Gamma \Delta u\,dx
= \int_{\partial\Omega}\left(\Gamma \frac{\partial u}{\partial \nu} - u \frac{\partial \Gamma}{\partial \nu}\right)ds
+ \int_{\partial B_\rho}\left(\Gamma \frac{\partial u}{\partial \nu} - u \frac{\partial \Gamma}{\partial \nu}\right)ds.
\end{equation}

Now

\[
\int_{\partial B_\rho} \Gamma \frac{\partial u}{\partial \nu}\,ds
= \Gamma(\rho)\int_{\partial B_\rho} \frac{\partial u}{\partial \nu}\,ds
\leq n\omega_n \rho^{n-1}\Gamma(\rho)\sup_{B_\rho}|Du|\to 0 \qquad \text{as }\rho\to 0
\]

and

\[
\int_{\partial B_\rho} u \frac{\partial \Gamma}{\partial \nu}\,ds
= -\Gamma'(\rho)\int_{\partial B_\rho} u\,ds \qquad (\text{recall that } \nu \text{ is outer normal to } \Omega-B_\rho)
\]
\[
= \frac{-1}{n\omega_n \rho^{n-1}}\int_{\partial B_\rho} u\,ds \to -u(y) \qquad \text{as }\rho\to 0.
\]

Hence letting $\rho$ tend to zero in (2.15) we arrive at Green's representation formula:

\begin{equation}
\tag{2.16}
u(y)= \int_{\partial\Omega}\left(u \frac{\partial \Gamma}{\partial \nu}(x-y)-\Gamma(x-y)\frac{\partial u}{\partial \nu}\right)ds + \int_{\Omega}\Gamma(x-y)\Delta u\,dx, \qquad (y\in \Omega).
\end{equation}

For an integrable function $f$, the integral $\int_\Omega \Gamma(x-y)f(x)\,dx$ is called the Newtonian

potential with density $f$. If $u$ has compact support in $\mathbb{R}^n$, then (2.16) yields the
frequently useful representation formula,

\begin{equation}
\tag{2.17}
u(y)= \int_{\Omega}\Gamma(x-y)\Delta u(x)\,dx.
\end{equation}

For harmonic $u$, we also obtain the representation

\begin{equation}
\tag{2.18}
u(y)= \int_{\partial\Omega}\left(u \frac{\partial \Gamma}{\partial \nu}(x-y)-\Gamma(x-y)\frac{\partial u}{\partial \nu}\right)ds, \qquad (y\in \Omega).
\end{equation}

Since the integrand above is infinitely differentiable and, in fact, also analytic with
respect to $y$, it follows that $u$ is also analytic in $\Omega$. Thus harmonic functions are
analytic throughout their domain of definition and therefore uniquely determined
by their values in any open subset.

Now suppose that $h\in C^1(\overline{\Omega})\cap C^2(\Omega)$ satisfies $\Delta h=0$ in $\Omega$. Then again by
Green's second identity (2.11) we obtain

\begin{equation}
\tag{2.19}
-\int_{\partial\Omega}\left(u\frac{\partial h}{\partial \nu}-h\frac{\partial u}{\partial \nu}\right)ds
= \int_\Omega h\,\Delta u\,dx.
\end{equation}


% PDF page 31
\source{gilbarg-trudinger:page-0031}



Writing $G=\Gamma+h$ and adding (2.16) and (2.19) we then obtain a more general version of Green's representation formula:

\begin{equation}
u(y)=\int_{\partial \Omega}\left(u\,\frac{\partial G}{\partial \nu}-G\,\frac{\partial u}{\partial \nu}\right)\diff s+\int_{\Omega}G\,\Delta u\diff x.
\tag{2.20}
\end{equation}

If in addition $G=0$ on $\partial \Omega$ we have

\begin{equation}
u(y)=\int_{\partial \Omega}u\,\frac{\partial G}{\partial \nu}\diff s+\int_{\Omega}G\,\Delta u\diff x
\tag{2.21}
\end{equation}

and the function $G=G(x,y)$ is called the (Dirichlet) Green's function for the domain $\Omega$, sometimes also called the Green's function of the first kind for $\Omega$. By Theorem 2.4, the Green's function is unique and from the formula (2.21) its existence implies a representation for a $C^1(\overline{\Omega})\cap C^2(\Omega)$ harmonic function in terms of its boundary values.

2.5. The Poisson Integral

When the domain $\Omega$ is a ball the Green's function can be explicitly determined by the method of images and leads to the well known Poisson integral representation for harmonic functions in a ball. Namely, let $B_R=B_R(0)$ and for $x\in B_R$, $x\neq 0$ let

\begin{equation}
\ubar{x}=\frac{R^2}{|x|^2}x
\tag{2.22}
\end{equation}

denote its inverse point with respect to $B_R$; if $x=0$, take $\ubar{x}=\infty$. It is then easily verified that the Green's function for $B_R$ is given by

\begin{equation}
G(x,y)=
\begin{cases}
\Gamma(|x-y|)-\Gamma\!\left(\dfrac{|y|}{R}\,|x-\bar{y}|\right), & y\neq 0,\\[1.25ex]
\Gamma(|x|)-\Gamma(R), & y=0.
\end{cases}
\tag{2.23}
\end{equation}

\[
=\Gamma\!\left(\sqrt{|x|^2+|y|^2-2x\cdot y}\right)-\Gamma\!\left(\sqrt{\left(\dfrac{|x||y|}{R}\right)^2+R^2-2x\cdot y}\right)
\]
for all $x,y\in B_R$, $x\neq y$.

The function $G$ defined by (2.23) has the properties

\begin{equation}
G(x,y)=G(y,x), \qquad G(x,y)\leq 0 \quad \text{for } x,y\in \overline{B}_R.
\tag{2.24}
\end{equation}


% PDF page 32
\source{gilbarg-trudinger:page-0032}



Furthermore, direct calculation shows that at $x \in \partial B_R$ the normal derivative of $G$ is given by

\begin{equation}
\frac{\partial G}{\partial \nu}=\frac{\partial G}{\partial |x|}=\frac{R^2-|y|^2}{n\omega_n R}|x-y|^{-n}\geq 0.
\tag{2.25}
\end{equation}

Hence if $u \in C^2(B_R)\cap C^1(\overline{B_R})$ is harmonic, we have by (2.21) the \textit{Poisson integral formula}:

\begin{equation}
u(y)=\frac{R^2-|y|^2}{n\omega_n R}\int_{\partial B_R}\frac{u\,d s_x}{|x-y|^n}
\tag{2.26}
\end{equation}

The right hand side of formula (2.26) is called the Poisson integral of $u$. A simple approximation argument shows that the Poisson integral formula continues to hold for $u \in C^2(B_R)\cap C^0(\overline{B_R})$. Note that by taking $y=0$, we recover the mean value theorem for harmonic functions. In fact all the previous theorems of this chapter could have been derived as consequences of the representation (2.21) with $\Omega = B_R(0)$.

To establish the existence of solutions of the classical Dirichlet problem for balls we need the converse result to the representation (2.26), and we prove this now.

\begin{theorem}\label{gt.thm2.6}\dcref{Theorem 2.6}\group{chapter-2}\source{gilbarg-trudinger:page-0032}
\textit{Let $B=B_R(0)$ and $\varphi$ be a continuous function on $\partial B$. Then the function $u$ defined by}

\begin{equation}
u(x)=
\begin{cases}
\dfrac{R^2-|x|^2}{n\omega_n R}\int_{\partial B}\dfrac{\varphi(y)\,d s_y}{|x-y|^n} & \text{for } x \in B \\[1.25ex]
\varphi(x) & \text{for } x \in \partial B
\end{cases}
\tag{2.27}
\end{equation}

\textit{belongs to $C^2(B)\cap C^0(\overline{B})$ and satisfies $\Delta u=0$ in $B$.}

\begin{proof}\uses{gt.thm2.5}
That $u$ defined by (2.27) is harmonic in $B$ is evident from the fact that $G$, and hence $\partial G/\partial \nu$, is harmonic in $x$, or it may be verified by direct calculation. To establish the continuity of $u$ on $\partial B$, we use the Poisson formula (2.26) for the special case $u=1$ to obtain the identity.

\begin{equation}
\int_{\partial B}K(x,y)\,d s_y=1 \qquad \text{for all } x \in B
\tag{2.28}
\end{equation}

where $K$ is the \textit{Poisson kernel}

\begin{equation}
K(x,y)=\frac{R^2-|x|^2}{n\omega_n R|x-y|^n}; \quad x \in B,\ y \in \partial B.
\tag{2.29}
\end{equation}

Of course the integral in (2.28) may be evaluated directly but this is a complicated calculation. Now let $x_0 \in \partial B$ and $\varepsilon$ be an arbitrary positive number. Choose $\delta>0$


% PDF page 33
\source{gilbarg-trudinger:page-0033}

so that $|\varphi(x)-\varphi(x_0)|<\varepsilon$ if $|x-x_0|<\delta$ and let $|\varphi|\le M$ on $\partial B$. Then if $|x-x_0|<\delta/2$, we have by (2.27) and (2.28)
\[
\begin{aligned}
|u(x)-u(x_0)|
&= \left|\int_{\partial B} K(x,y)(\varphi(y)-\varphi(x_0))\,ds_y\right| \\
&\le \int_{|y-x_0|\le\delta} K(x,y)|\varphi(y)-\varphi(x_0)|\,ds_y \\
&\quad + \int_{|y-x_0|>\delta} K(x,y)|\varphi(y)-\varphi(x_0)|\,ds_y \\
&\le \varepsilon + \frac{2M(R^2-|x|^2)R^{n-2}}{(\delta/2)^n}.
\end{aligned}
\]
If now $|x-x_0|$ is sufficiently small it is clear that $|u(x)-u(x_0)|<2\varepsilon$ and hence $u$ is continuous at $x_0$. Consequently $u\in C^0(\overline{B})$ as required. $\square$
\end{proof}
\end{theorem}

We note that the preceding argument is local; that is, if $\varphi$ is only bounded and integrable on $\partial B$, and continuous at $x_0$, then $u(x)\to\varphi(x_0)$ as $x\to x_0$.

2.6. Convergence Theorems

We consider now some immediate consequences of the Poisson integral formula. The following three theorems will not however be required for the later development. We show first that harmonic functions can in fact be characterized by their mean value property.

\begin{theorem}\label{gt.thm2.7}\dcref{Theorem 2.7}\group{chapter-2}\source{gilbarg-trudinger:page-0033}
$A\ C^0(\Omega)$ function $u$ is harmonic if and only if for every ball $B=B_R(y)\Subset\Omega$ it satisfies the mean value property,
\[
(2.30)\qquad u(y)=\frac{1}{n\omega_nR^{n-1}}\int_{\partial B} u\,ds.
\]

\begin{proof}\uses{gt.thm2.6, gt.thm2.2, gt.thm2.3, gt.thm2.4}
By Theorem 2.6, there exists for any ball $B\Subset\Omega$ a harmonic function $h$ such that $h=u$ on $\partial B$. The difference $w=u-h$ will then be a function satisfying the mean value property on any ball in $B$. Consequently the maximum principle and uniqueness results of Theorems 2.2, 2.3 and 2.4 apply to $w$ since the mean value inequalities were the only properties of harmonic functions used in their derivation. Hence $w=0$ in $B$ and consequently $u$ must be harmonic in $\Omega$. $\square$
\end{proof}
\end{theorem}

As an immediate consequence of the preceding theorem we have:

\begin{theorem}\label{gt.thm2.8}\dcref{Theorem 2.8}\group{chapter-2}\source{gilbarg-trudinger:page-0033}
The limit of a uniformly convergent sequence of harmonic functions is harmonic.
\end{theorem}

% PDF page 34
\source{gilbarg-trudinger:page-0034}


It follows from Theorem 2.8, that if $\{u_n\}$ is a sequence of harmonic functions in a bounded domain $\Omega$, with continuous boundary values $\{\varphi_n\}$ which converge uniformly on $\partial\Omega$ to a function $\varphi$, then the sequence $\{u_n\}$ converges uniformly (by the maximum principle) to a harmonic function having the boundary values $\varphi$ on $\partial\Omega$. By means of Harnack's inequality, Theorem 2.5, we can also derive, from Theorem 2.8, \textit{Harnack's convergence theorem}.

\medskip

\begin{theorem}\label{gt.thm2.9}\dcref{Theorem 2.9}\group{chapter-2}\source{gilbarg-trudinger:page-0034}
\textit{Let $\{u_n\}$ be a monotone increasing sequence of harmonic functions in a domain $\Omega$ and suppose that for some point $y \in \Omega$, the sequence $\{u_n(y)\}$ is bounded. Then the sequence converges uniformly on any bounded subdomain $\Omega' \Subset \Omega$ to a harmonic function.}

\medskip

\begin{proof}\uses{gt.thm2.5, gt.thm2.8}
The sequence $\{u_n(y)\}$ will converge, so that for arbitrary $\varepsilon>0$ there is a number $N$ such that $0 \leq u_m(y)-u_n(y)<\varepsilon$ for all $m \geq n > N$. But then by Theorem 2.5, we must have

\[
\sup_{\Omega'} |u_m(x)-u_n(x)| < C\varepsilon
\]

for some constant $C$ depending on $\Omega'$ and $\Omega$. Consequently $\{u_n\}$ converges uniformly and by virtue of Theorem 2.8, the limit function is harmonic. \quad $\square$

\bigskip

\noindent\textbf{2.7.}\ \textbf{Interior Estimates of Derivatives}

\medskip

By direct differentiation of the Poisson integral it is possible to obtain interior derivative estimates for harmonic functions. Alternatively, such estimates also follow from the mean value theorem. For let $u$ be harmonic in $\Omega$ and $B=B_R(y)\subset\subset \Omega$. Since the gradient $Du$ is also harmonic in $\Omega$ it follows by the mean value and divergence theorems that

\[
Du(y)=\frac{1}{\omega_n R^n}\int_B Du\,dx=\frac{1}{\omega_n R^n}\int_{\partial B} uw\,ds,
\]

\[
|Du(y)|\leq \frac{n}{R}\sup_{\partial B}|u|
\]

and hence

\[
(2.31)\qquad |Du(y)|\leq \frac{n}{d_y}\sup_{\Omega}|u|,
\]

where $d_y=\operatorname{dist}(y,\partial\Omega)$. By successive application of the estimate (2.31) in equally spaced nested balls we obtain an estimate for higher order derivatives:
\end{proof}
\end{theorem}


% PDF page 35
\source{gilbarg-trudinger:page-0035}

2.8. The Dirichlet Problem; the Method of Subharmonic Functions

We are in a position now to approach the question of existence of solutions of the classical Dirichlet problem in arbitrary bounded domains. The treatment here will be accomplished by Perron's method of subharmonic functions [PE] which relies heavily on the maximum principle and the solvability of the Dirichlet problem in balls. The method has a number of attractive features in that it is elementary, it separates the interior existence problem from that of the boundary behaviour of solutions, and it is easily extended to more general classes of second order elliptic equations. There are other well known approaches to existence theorems such as the method of integral equations, treated for example in the books [KE 2] [GU], and the variational or Hilbert space approach which we describe in a more general context in Chapter 8.

The definition of $C^2(\Omega)$ subharmonic and superharmonic function is generalized as follows. A $C^0(\Omega)$ function $u$ will be called subharmonic (superharmonic) in $\Omega$ if for every ball $B \subset \subset \Omega$ and every function $h$ harmonic in $B$ satisfying $u \leq (\geq) h$ on $\partial B$, we also have $u \leq (\geq) h$ in $B$. The following properties of $C^0(\Omega)$ subharmonic functions are readily established:

\quad (i) If $u$ is subharmonic in a domain $\Omega$, it satisfies the strong maximum principle in $\Omega$; and if $v$ is superharmonic in a bounded domain $\Omega$ with $v \geq u$ on $\partial \Omega$, then either $v > u$ throughout $\Omega$ or $v \equiv u$. To prove the latter assertion, suppose the contrary. Then at some point $x_0 \in \Omega$ we have

\[
(u-v)(x_0)=\sup_{\Omega}(u-v)=M\geq 0,
\]

% PDF page 36
\source{gilbarg-trudinger:page-0036}

and we may assume there is a ball $B=B(x_0)$ such that $u-v\not\equiv M$ on $\partial B$. Letting $\bar u$, $\bar v$ denote the harmonic functions respectively equal to $u$, $v$ on $\partial B$ (Theorem 2.6), one sees that

\[
M \geq \sup_{\partial B}(\bar u-\bar v) \geq (\bar u-\bar v)(x_0) \geq (u-v)(x_0)=M,
\]

and hence the equality holds throughout. By the strong maximum principle for harmonic functions (Theorem 2.2) it follows that $\bar u-\bar v\equiv M$ in $B$ and hence $u-v\equiv M$ on $\partial B$, which contradicts the choice of $B$.

\noindent (ii) Let $u$ be subharmonic in $\Omega$ and $B$ be a ball strictly contained in $\Omega$. Denote by $\bar u$ the harmonic function in $B$ (given by the Poisson integral of $u$ on $\partial B$) satisfying $\bar u=u$ on $\partial B$. We define in $\Omega$ the harmonic lifting of $u$ (in $B$) by

\[
U(x)=\begin{cases}
\bar u(x), & x\in B \\
 u(x), & x\in \Omega-B.
\end{cases}
\tag{2.33}
\]

Then the function $U$ is also subharmonic in $\Omega$. For consider an arbitrary ball $B'\subset\subset\Omega$ and let $h$ be a harmonic function in $B'$ satisfying $h\geq U$ on $\partial B'$. Since $u\leq U$ in $B'$ we have $u\leq h$ in $B'$ and hence $U\leq h$ in $B'-B$. Also since $U$ is harmonic in $B$, we have by the maximum principle $U\leq h$ in $B\cap B'$. Consequently $U\leq h$ in $B'$ and $U$ is subharmonic in $\Omega$.

\noindent (iii) Let $u_1, u_2, \dots, u_N$ be subharmonic in $\Omega$. Then the function $u(x)=\max\{u_1(x),\dots,u_N(x)\}$ is also subharmonic in $\Omega$. This is a trivial consequence of the definition of subharmonicity. Corresponding results for superharmonic functions are obtained by replacing $u$ by $-u$ in properties (i), (ii) and (iii).

Now let $\Omega$ be bounded and $\varphi$ be a bounded function on $\partial\Omega$. A $C^0(\overline\Omega)$ subharmonic function $u$ is called a subfunction relative to $\varphi$ if it satisfies $u\leq \varphi$ on $\partial\Omega$. Similarly a $C^0(\overline\Omega)$ superharmonic function is called a superfunction relative to $\varphi$ if it satisfies $u\geq \varphi$ on $\partial\Omega$. By the maximum principle every subfunction is less than or equal to every superfunction. In particular, constant functions $\leq \inf_{\partial\Omega}\varphi$ ($\geq \sup_{\partial\Omega}\varphi$) are subfunctions (superfunctions). Let $S_\varphi$ denote the set of subfunctions relative to $\varphi$. The basic result of the Perron method is contained in the following theorem.

\medskip
\begin{theorem}\label{gt.thm2.12}\dcref{Theorem 2.12}\group{chapter-2}\source{gilbarg-trudinger:page-0036}
\emph{The function $u(x)=\sup_{v\in S_\varphi}v(x)$ is harmonic in $\Omega$.}

\medskip
\begin{proof}\uses{gt.thm2.9}
By the maximum principle any function $v\in S_\varphi$ satisfies $v\leq \sup \varphi$, so that $u$ is well defined. Let $y$ be an arbitrary fixed point of $\Omega$. By the definition of $u$, there exists a sequence $\{v_n\}\subset S_\varphi$ such that $v_n(y)\to u(y)$. By replacing $v_n$ with $\max\,(v_n,\inf\varphi)$, we may assume that the sequence $\{v_n\}$ is bounded. Now choose $R$ so that the ball $B=B_R(y)\subset\subset\Omega$ and define $V_n$ to be the harmonic lifting of $v_n$ in $B$ according to (2.33). Then $V_n\in S_\varphi$, $V_n(y)\to u(y)$ and by Theorem 2.11 the sequence $\{V_n\}$ contains a subsequence $\{V_{n_k}\}$ converging uniformly in any ball $B_\rho(y)$ with $\rho<R$ to a function $v$ that is harmonic in $B$. Clearly $v\leq u$ in $B$ and $v(y)=u(y)$. We claim now that in fact $v=u$ in $B$. For suppose $v(z)<u(z)$ at some $z\in B$. Then there exists


% PDF page 37
\source{gilbarg-trudinger:page-0037}

\begin{quote}
a function $\bar{u} \in S_{\varphi}$ such that $v(z) < \bar{u}(z)$. Defining $w_k = \max(\bar{u}, v_{n_k})$ and also the harmonic liftings $W_k$ as in (2.33), we obtain as before a subsequence of the sequence $\{W_k\}$ converging to a harmonic function $w$ satisfying $v \leq w \leq u$ in $B$ and $v(y)=w(y)=u(y)$. But then by the maximum principle we must have $v=w$ in $B$. This contradicts the definition of $\bar{u}$ and hence $u$ is harmonic in $\Omega$. $\square$
\end{quote}
\end{proof}
\end{theorem}

The preceding result exhibits a harmonic function which is a prospective solution (called the Perron solution) of the classical Dirichlet problem: $\Delta u = 0$, $u = \varphi$ on $\partial \Omega$. Indeed, if the Dirichlet problem is solvable, its solution is identical with the Perron solution. For let $w$ be the presumed solution. Then clearly $w \in S_{\varphi}$ and by the maximum principle $w \geq u$ for all $u \in S_{\varphi}$. We note here also that the proof of Theorem 2.12 could have been based on the Harnack convergence theorem, Theorem 2.9, instead of the compactness theorem, Theorem 2.11; (see Problem 2.10).

In the Perron method the study of boundary behaviour of the solution is essentially separate from the existence problem. The continuous assumption of boundary values is connected to the geometric properties of the boundary through the concept of barrier function. Let $\xi$ be a point of $\partial \Omega$. Then a $C^0(\bar{\Omega})$ function $w = w_{\xi}$ is called a barrier at $\xi$ relative to $\Omega$ if:

\begin{enumerate}
\item[(i)] $w$ is superharmonic in $\Omega$;
\item[(ii)] $w > 0$ in $\bar{\Omega} - \xi$; $w(\xi)=0$.
\end{enumerate}

A more general definition of barrier requires only that the superharmonic function $w$ be continuous and positive in $\Omega$, and that $w(x) \to 0$ as $x \to \xi$. The results of this section are valid for these weak barriers as well (see [HL, p. 168], for example). An important feature of the barrier concept is that it is a local property of the boundary $\partial \Omega$. Namely, let us define $w$ to be a local barrier at $\xi \in \partial \Omega$ if there is a neighborhood $N$ of $\xi$ such that $w$ satisfies the above definition in $\Omega \cap N$. Then a barrier at $\xi$ relative to $\Omega$ can be defined as follows. Let $B$ be a ball satisfying $\xi \in B \subset N$ and $m = \inf_{N-B} w > 0$. The function

\[
\bar{w}(x) =
\begin{cases}
\min(m, w(x)), & x \in \bar{\Omega} \cap B \\
m, & x \in \bar{\Omega} - B
\end{cases}
\]

is then a barrier at $\xi$ relative to $\Omega$, as one sees by confirming properties (i) and (ii). Indeed, $\bar{w}$ is continuous in $\bar{\Omega}$ and is superharmonic in $\Omega$ by property (iii) of subharmonic functions; property (ii) is immediate.

A boundary point will be called regular (with respect to the Laplacian) if there exists a barrier at that point.

The connection between the barrier and boundary behavior of solutions is contained in the following.

\medskip

\begin{lemma}\label{gt.lem2.13}\dcref{Lemma 2.13}\group{chapter-2}\source{gilbarg-trudinger:page-0037}
\textit{Let $u$ be the harmonic function defined in $\Omega$ by the Perron method (Theorem 2.12). If $\xi$ is a regular boundary point of $\Omega$ and $\varphi$ is continuous at $\xi$, then $u(x) \to \varphi(\xi)$ as $x \to \xi$.}

% PDF page 38
\source{gilbarg-trudinger:page-0038}
\begin{quote}
    \begin{proof}\uses{gt.thm2.12}
Choose $\varepsilon>0$, and let $M=\sup |\varphi|$. Since $\zeta$ is a regular boundary point, there is a barrier $w$ at $\zeta$ and, by virtue of the continuity of $\varphi$, there are constants $\delta$ and $k$ such that $|\varphi(x)-\varphi(\zeta)|<\varepsilon$ if $|x-\zeta|<\delta$, and $kw(x)\geq 2M$ if $|x-\zeta|\geq \delta$. The functions $\varphi(\zeta)+\varepsilon+kw$, $\varphi(\zeta)-\varepsilon-kw$ are respectively superfunction and subfunction relative to $\varphi$. Hence from the definition of $u$ and the fact that every superfunction dominates every subfunction, we have in $\Omega$.
\[
\varphi(\zeta)-\varepsilon-kw(x)\leq u(x)\leq \varphi(\zeta)+\varepsilon+kw(x)
\]
or
\[
|u(x)-\varphi(\zeta)|\leq \varepsilon+kw(x).
\]

Since $w(x)\to 0$ as $x\to \zeta$, we obtain $u(x)\to \varphi(\zeta)$ as $x\to \zeta$. $\square$
\end{proof}
This leads immediately to
\end{quote}
\end{lemma}

\begin{theorem}\label{gt.thm2.14}\dcref{Theorem 2.14}\group{chapter-2}\source{gilbarg-trudinger:page-0038}
\emph{The classical Dirichlet problem in a bounded domain is solvable for arbitrary continuous boundary values if and only if the boundary points are all regular.}

\begin{quote}
\begin{proof}\uses{gt.lem2.13}
If the boundary values $\varphi$ are continuous and the boundary $\partial\Omega$ consists of regular points, the preceding lemma states that the harmonic function provided by the Perron method solves the Dirichlet problem. Conversely, suppose that the Dirichlet problem is solvable for all continuous boundary values. Let $\zeta\in\partial\Omega$. Then the function $\varphi(x)=|x-\zeta|$ is continuous on $\partial\Omega$ and the harmonic function solving the Dirichlet problem in $\Omega$ with boundary values $\varphi$ is obviously a barrier at $\zeta$. Hence $\zeta$ is regular, as are all points of $\partial\Omega$. $\square$
\end{proof}
\end{quote}
\end{theorem}

\noindent The important question remains: For what domains are the boundary points regular? It turns out that general sufficient conditions can be stated in terms of local geometric properties of the boundary. We mention some of these conditions below.

\noindent If $n=2$, consider a boundary point $z_{0}$ of a bounded domain $\Omega$ and take the origin at $z_{0}$ with polar coordinates $r,\theta$. Suppose there is a neighborhood $N$ of $z_{0}$ such that a single valued branch of $\theta$ is defined in $\Omega\cap N$, or in a component of $\Omega\cap N$ having $z_{0}$ on its boundary. One sees that
\[
w=-\Re \frac{1}{\log z}=-\frac{\log r}{\log^{2} r+\theta^{2}}
\]
is a (weak) local barrier at $z_{0}$ and hence $z_{0}$ is a regular point. In particular, $z_{0}$ is a regular boundary point if it is the endpoint of a simple arc lying in the exterior of $\Omega$. Thus the Dirichlet problem in the plane is always solvable for continuous boundary values in a (bounded) domain whose boundary points are each accessible from the exterior by a simple arc. More generally, the same barrier shows that the boundary value problem is solvable if every component of the complement of the domain consists of more than a single point. Examples of such domains are domains bounded by a finite number of simple closed curves. Another is the unit disc slit along an arc; in this case the boundary values can be assigned on opposite sides of the slit.

% PDF page 39
\source{gilbarg-trudinger:page-0039}

For higher dimensions the situation is substantially different and the Dirichlet
problem cannot be solved in corresponding generality. Thus, an example due to
Lebesgue shows that a closed surface in three dimensions with a sufficiently sharp
inward directed cusp has a non-regular point at the tip of the cusp; (see for
example [CH]).

A simple sufficient condition for solvability in a bounded domain $\Omega \subset \mathbb{R}^n$ is
that $\Omega$ satisfy the \emph{exterior sphere condition}; that is, for every point $\xi \in \partial\Omega$, there
exists a ball $B = B_R(y)$ satisfying $\bar B \cap \bar\Omega = \xi$. If such a condition is fulfilled, then the
function $w$ given by

\begin{equation}
\tag{2.34}
w(x)=
\begin{cases}
R^{2-n}-|x-y|^{2-n} & \text{for } n \geq 3 \\
\log \dfrac{|x-y|}{R} & \text{for } n=2
\end{cases}
\end{equation}

will be a barrier at $\xi$. Consequently the boundary points of a domain with $C^2$
boundary are all regular points; (see Problem 2.11).

\section*{2.9. Capacity}

The physical concept of capacity provides another means of characterizing regular
and exceptional boundary points. Let $\Omega$ be a bounded domain in $\mathbb{R}^n(n \geq 3)$ with
smooth boundary $\partial\Omega$, and let $u$ be the harmonic function (often called the con-
ductor potential) defined in the complement of $\bar\Omega$ and satisfying the boundary
conditions $u=1$ on $\partial\Omega$ and $u=0$ at infinity. The existence of $u$ is easily established
as the (unique) limit of harmonic functions $u'$ in an expanding sequence of bounded
domains having $\partial\Omega$ as an inner boundary (on which $u' = 1$) and with outer
boundaries (on which $u' = 0$) tending to infinity. If $\Sigma$ denotes $\partial\Omega$ or any smooth
closed surface enclosing $\Omega$, then the quantity

\begin{equation}
\tag{2.35}
\operatorname{cap}\Omega = - \int_{\Sigma} \frac{\partial u}{\partial \nu}\, ds
= \int_{\mathbb{R}^n \setminus \Omega} |Du|^2\, dx
\qquad \nu = \text{outer normal}
\end{equation}

is defined to be the capacity of $\Omega$. In electrostatics, $\operatorname{cap}\Omega$ is within a constant factor
the total electric charge on the conductor $\partial\Omega$ held at unit potential (relative to
infinity).

Capacity can also be defined for domains with nonsmooth boundaries and for
any compact set as the (unique) limit of the capacities of a nested sequence of
approximating smoothly bounded domains. Equivalent definitions of capacity can
be given directly without use of approximating domains (e.g., see [LK]). In
particular, we have the variational characterization

\begin{equation}
\tag{2.36}
\operatorname{cap}\Omega = \inf_{v \in \mathcal{K}} \int |Dv|^2.
\end{equation}

% PDF page 40
\source{gilbarg-trudinger:page-0040}

\[
K = \{v \in C_0^1(\mathbb{R}^n)\mid v = 1 \text{ on } \Omega\}.
\]

To investigate the regularity of a point \(x_0 \in \partial\Omega\), consider for any fixed \(\lambda \in (0,1)\)
the capacity

\[
C_j = \operatorname{cap}\{x \notin \Omega \mid |x - x_0| \le \lambda^j\}.
\]

The Wiener criterion states that \(x_0\) is a regular boundary point of \(\Omega\) if and only if the
series

\[
(2.37)\qquad \sum_{j=0}^{\infty} \frac{C_j}{\lambda^{j(n-2)}}
\]

diverges.

For a discussion of capacity and proof of the Wiener criterion we refer to the
literature, e.g., [KE 2, LK]. In Chapter 8 this condition for regularity will be proved
for a general class of elliptic operators in divergence form.

\paragraph{Problems}

\textbf{2.1.} Derive the weak maximum principle for \(C^2(\Omega)\) subharmonic functions from
a consideration of necessary conditions for a relative maximum.

\textbf{2.2.} Prove that if \(\Delta u = 0\) in \(\Omega\) and \(u = \partial u/\partial \nu = 0\) on an open, smooth portion of \(\partial\Omega\),
then \(u\) is identically zero.

\textbf{2.3.} Let \(G\) be the Green's function for a bounded domain \(\Omega\). Prove

\begin{enumerate}
\item[(a)] \(G(x,y)=G(y,x)\) for all \(x,y\in\Omega,\ x\ne y;\)
\item[(b)] \(G(x,y)<0\) for all \(x,y\in\Omega,\ x\ne y;\)
\item[(c)] \(\displaystyle \int_{\Omega} G(x,y)f(y)\,dy \to 0\) as \(x\to \partial\Omega\), if \(f\) is bounded and integrable on \(\Omega\).
\end{enumerate}

\textbf{2.4.} (\textit{Schwarz reflection principle.}) Let \(\Omega^+\) be a subdomain of the half-space
\(x_n>0\) having as part of its boundary an open section \(T\) of the hyperplane \(x_n=0\).
Suppose that \(u\) is harmonic in \(\Omega^+\), continuous in \(\Omega^+ \cup T\), and that \(u=0\) on \(T\).
Show that the function \(U\) defined by

\[
U(x_1,\dots,x_n)=
\begin{cases}
u(x_1,\dots,x_n), & x_n\ge 0\\
-u(x_1,\dots,-x_n), & x_n<0
\end{cases}
\]

% PDF page 41
\source{gilbarg-trudinger:page-0041}

is harmonic in the domain $\Omega^+ \cup T \cup \Omega^-$, where $\Omega^-$ is the reflection of $\Omega^+$ in $x_n=0$ (i.e., $\Omega^-=\{(x_1,\ldots,x_n)\in\mathbb{R}^n\mid (x_1,\ldots,-x_n)\in\Omega^+\}$).

\begin{enumerate}
\setcounter{enumi}{4}
\item[2.5.] Determine the Green's function for the annular region bounded by two concentric spheres in $\mathbb{R}^n$.

\item[2.6.] Let $u$ be a non-negative harmonic function in a ball $B_R(0)$. Deduce from the Poisson integral formula, the following version of Harnack's inequality

\[
\frac{R^{n-2}(R-|x|)}{(R+|x|)^{n-1}}\,u(0)\le u(x)\le \frac{R^{n-2}(R+|x|)}{(R-|x|)^{n-1}}\,u(0).
\]

\item[2.7.] Show that a $C^0(\Omega)$ function $u$ is subharmonic in $\Omega$ if and only if it satisfies the mean value inequality locally; that is, for every $y\in\Omega$ there exists $\delta=\delta(y)>0$ such that

\[
u(y)\le \frac{1}{n\omega_n R^{n-1}}\int_{\partial B_R(y)} u\,ds \qquad \text{for all } R\le \delta.
\]

\item[2.8.] An integrable function $u$ in a domain $\Omega$ is called weakly harmonic (subharmonic, superharmonic) in $\Omega$ if

\[
\int_\Omega u\,\Delta\varphi\,dx \equiv (\ge,\le) 0
\]

for all functions $\varphi\ge 0$ in $C^2(\Omega)$ having compact support in $\Omega$. Show that a $C^0(\Omega)$ weakly harmonic (subharmonic, superharmonic) function is harmonic (subharmonic, superharmonic).

\item[2.9.] Show that for $C^2(\Omega)$ functions $u$, the conditions: (i) $\Delta u\ge 0$ in $\Omega$; (ii) $u$ is subharmonic in $\Omega$; (iii) $u$ is weakly subharmonic in $\Omega$, are equivalent.

\item[2.10.] Prove Theorem 2.12 using Theorem 2.9 instead of Theorem 2.11.

\item[2.11.] Show that a domain $\Omega$ with $C^2$ boundary $\partial\Omega$ satisfies an exterior sphere condition.

\item[2.12.] Show that the Dirichlet problem is solvable for any domain $\Omega$ satisfying an exterior cone condition; that is, for every point $\xi\in\partial\Omega$ there exists a finite right circular cone $K$, with vertex $\xi$, satisfying $\overline{K}\cap\overline{\Omega}=\xi$. At each point $\xi\in\partial\Omega$ taken as origin, show that a suitable local barrier can be chosen in the form $w=r^\alpha f(\theta)$ where $\theta$ is the polar angle.

\item[2.13.] Let $u$ be harmonic in $\Omega\subset\mathbb{R}^n$. Use the argument leading to (2.31) to prove the interior gradient bound,

\[
|Du(x_0)|\le \frac{n}{d_0}\left[\sup_{\Omega} u-u(x_0)\right],\qquad d_0=\operatorname{dist}(x_0,\partial\Omega).
\]
\end{enumerate}

% PDF page 42
\source{gilbarg-trudinger:page-0042}

If $u \geq 0$ in $\Omega$ infer that

\[
\lvert Du(x_0)\rvert \leq \frac{n}{d_0}u(x_0).
\]

\paragraph{2.14.} (a) Prove \emph{Liouville's theorem}: A harmonic function defined over $\mathbb{R}^n$ and bounded above is constant.

(b) If $n=2$ prove that the Liouville theorem in part (a) is valid for \emph{subharmonic} functions.

(c) If $n>2$ show that a bounded subharmonic function defined over $\mathbb{R}^n$ need not be constant.

\paragraph{2.15.} Let $u \in C^2(\overline{\Omega})$, $u = 0$ on $\partial \Omega \in C^1$. Prove the interpolation inequality: For every $\varepsilon > 0$,

\[
\int_{\Omega} \lvert Du\rvert^2\,dx \leq \varepsilon \int_{\Omega} (\Delta u)^2\,dx + \frac{1}{4\varepsilon}\int_{\Omega} u^2\,dx.
\]

\paragraph{2.16.} Prove Theorem 2.12 by finding in every ball $B \Subset \Omega$ a monotone increasing sequence of harmonic functions that are restrictions of subfunctions on $B$ and that converge uniformly to $u$ on a dense set of points in $B$. Hence show that Theorems 2.12 and 2.14 can be proved without use of the strong maximum principle.

\paragraph{2.17.} Show that the volume integral in (2.35) is defined, and prove the equivalence of the capacity definitions (2.35) and (2.36).

\paragraph{2.18.} Let $u$ be harmonic in (open, connected) $\Omega \subset \mathbb{R}^n$, and suppose $B_c(x_0) \subset \subset \Omega$.
If $a \leq b \leq c$, where $b^2 = ac$, show that

\[
\int_{\lvert \omega\rvert = 1} u(x_0 + a\omega)u(x_0 + c\omega)\,d\omega
=
\int_{\lvert \omega\rvert = 1} u^2(x_0 + b\omega)\,d\omega.
\]

Hence, conclude that if $u$ is constant in a neighbourhood it is identically constant.
(Cf. [GN].)

% PDF page 43
\source{gilbarg-trudinger:page-0043}
\chapter{The Classical Maximum Principle}

The purpose of this chapter is to extend the classical maximum principles for the
Laplace operator, derived in Chapter 2, to linear elliptic differential operators of
the form

\begin{equation}\tag{3.1}
Lu=a^{ij}(x)D_{ij}u+b^i(x)D_i u+c(x)u,\qquad a^{ij}=a^{ji},
\end{equation}
where $x=(x_1,\dots,x_n)$ lies in a domain $\Omega$ of $\mathbb{R}^n$, $n\ge 2$. It will be assumed, unless
otherwise stated, that $u$ belongs to $C^2(\Omega)$. The summation convention that repeated
indices indicate summation from $1$ to $n$ is followed here as it will be throughout.
 $L$ will always denote the operator (3.1).

We adopt the following definitions:

$L$ is elliptic at a point $x\in\Omega$ if the coefficient matrix $[a^{ij}(x)]$ is positive; that is, if
$\lambda(x),\Lambda(x)$ denote respectively the minimum and maximum eigenvalues of $[a^{ij}(x)]$,
then

\begin{equation}\tag{3.2}
0<\lambda(x)|\xi|^2\le a^{ij}(x)\xi_i\xi_j\le \Lambda(x)|\xi|^2
\end{equation}

for all $\xi=(\xi_1,\dots,\xi_n)\in\mathbb{R}^n-\{0\}$. If $\lambda>0$ in $\Omega$, then $L$ is elliptic in $\Omega$, and strictly
elliptic if $\lambda\ge \lambda_0>0$ for some constant $\lambda_0$. If $\Lambda/\lambda$ is bounded in $\Omega$, we shall call $L$
uniformly elliptic in $\Omega$. Thus the operator $D_{11}+x_1D_{22}$ is elliptic but not uniformly
elliptic in the half plane, $x_1>0$, while it is uniformly elliptic in strips of the form
$(\alpha,\beta)\times\mathbb{R}$ where $0<\alpha<\beta<\infty$.

Most results concerning elliptic operators of the form (3.1) require additional
conditions limiting the relative importance of the lower order terms $b^iD_i u, cu$ with
respect to the principal term $a^{ij}D_{ij}u$. The condition

\begin{equation}\tag{3.3}
\frac{|b^i(x)|}{\lambda(x)}\le \mathrm{const}<\infty,\qquad i=1,\dots,n,\quad x\in\Omega
\end{equation}

will be assumed throughout this chapter. By then considering $L'=\lambda^{-1}L$ in place
of $L$ we can reduce to the case where $\lambda=1$ and the $b^i$ are bounded. If, in addition,
$L$ is uniformly elliptic, we can also take the $a^{ij}$ to be bounded. Note that if the
coefficients $a^{ij}, b^i$ are continuous in $\Omega$, then on any bounded subdomain $\Omega'\Subset\Omega$,
$L$ is uniformly elliptic and (3.3) holds. The coefficient $c$ will also be subject to
restrictions but these will vary and consequently be indicated in the appropriate
hypotheses.
